\clearpage

\section*{Lectures conservées et questions éditoriales non tranchées}

Ces notes présentent des alternatives qui ne sont pas insérées comme corrections nouvelles. Elles sont distinctes de l’apparat numéroté : une question sans différence entre les formes diplomatique et corrigée ne crée aucune entrée de variante. « Needs the scan » signifie que la lecture matérielle exige le scan, absent de cette construction. La couche diplomatique fournie reste le témoin de la lecture imprimée.

\subsection*{ED-G27-x-hat — 7I:123}

Lecture de la source conservée : \textbackslash{}mathfrak X\_\{(\textbackslash{}hat x)\} in the supplied source layers

Alternative : earlier transcription \textbackslash{}mathcal X\_\{(\textbackslash{}mathfrak x)\}

Conserver le chapeau des deux couches fournies. Le commentaire antérieur hésitait sur la forme physique du caractère surfrappé ; seul le scan permettrait de vérifier cette question paléographique.

Statut : needs the scan. Source : \nolinkurl{conventions/GLOBAL_ISSUES_SGA.md:27}

\subsection*{ED-G27-missing-reference — 7I:128}

Lecture de la source conservée : is smooth by and therefore

Alternative : is smooth by 4.13(a) and therefore

Conserver la référence laissée vide dans le texte imprimé. Plusieurs résultats de lissité peuvent intervenir ; aucun témoin fourni ne donne le numéro manquant. La proposition 4.13(a) reste une suggestion éditoriale.

Statut : retained; authorial reference unresolved. Source : \nolinkurl{conventions/GLOBAL_ISSUES_SGA.md:27}

\subsection*{ED-G32-run-in-headings — 7I:163,165 (and analogous numbered headings)}

Lecture de la source conservée : Attached layers use separate paragraphs after numbered headings

Alternative : possible printed run-in headings with following prose on the same paragraph

Conserver les alinéas des couches fournies. Les corrections de mise en ligne ont été retenues à faible confiance et les autres occurrences n’ont pas été vérifiées. Le scan est nécessaire pour établir les retraits et les titres courants.

Statut : needs the scan. Source : \nolinkurl{conventions/GLOBAL_ISSUES_SGA.md:32}

\subsection*{ED-G36-E-or-X — 7I:180}

Lecture de la source conservée : d\_1\textasciicircum{}*d\_0\textasciicircum{}*(E)\textbackslash{}simeq(d\_0d\_1)\textasciicircum{}*(E)

Alternative : d\_1\textasciicircum{}*d\_0\textasciicircum{}*(X)\textbackslash{}simeq(d\_0d\_1)\textasciicircum{}*(X)

Conserver E. L’extension considérée juste avant est X, mais l’isomorphisme de transitivité est valable pour une extension quelconque E ; une notation générique de l’auteur reste possible.

Statut : retained; no determinate emendation. Source : \nolinkurl{conventions/GLOBAL_ISSUES_SGA.md:36}

\subsection*{ED-G36-psi-footnote — 7I:180}

Lecture de la source conservée : \textbackslash{}psi\_\{L.\}(A)

Alternative : \textbackslash{}underline\textbackslash{}Psi\_\{L.\}(A)

Conserver le psi minuscule de la couche diplomatique. Le texte principal emploie un Psi majuscule souligné, mais le dossier ne permet pas de distinguer une variante de glyphe d’une coquille. Le scan est nécessaire.

Statut : needs the scan. Source : \nolinkurl{conventions/GLOBAL_ISSUES_SGA.md:36}

\subsection*{ED-G37-duplicate-number — 7I:183}

Lecture de la source conservée : two displays tagged (3.2.3)

Alternative : second display (3.2.3\textbackslash{}prime) or another new tag

Conserver les deux numéros (3.2.3). Le numéro (3.2.4) est déjà employé ensuite ; aucun numéro voulu par l’auteur ne peut être rétabli avec certitude. Le diagramme et le morphisme c(X,alpha) sont distingués par leur contenu.

Statut : retained; intended numbering unresolved. Source : \nolinkurl{conventions/GLOBAL_ISSUES_SGA.md:37}

\subsection*{ED-G38-speck — 7II:202}

Lecture de la source conservée : (\textbackslash{}mathrm\{Sw\}\_\{I\textbackslash{}prime\},s\_\{I/H\}) without an extra lower dot

Alternative : a possible printed dot below the opening parenthesis

Ne pas ajouter de point sous la parenthèse. Aucune opération mathématique n’est définie par cette marque possible ; seul le scan peut distinguer ponctuation, encre et défaut matériel.

Statut : needs the scan. Source : \nolinkurl{conventions/GLOBAL_ISSUES_SGA.md:38}

\subsection*{ED-G41-exton-sion — 7I:195}

Lecture de la source conservée : extension (both attached layers)

Alternative : exton-sion (held low-confidence visual reading)

Conserver extension, lecture commune des deux couches fournies. La lecture exton-sion a été retenue à faible confiance ; le scan permettrait de distinguer un o d’un e endommagé et de vérifier le trait de coupure de fin de ligne.

Statut : needs the scan. Source : \nolinkurl{conventions/GLOBAL_ISSUES_SGA.md:41}

\subsection*{ED-G45-mathematical-repair — 7II:237}

Lecture de la source conservée : the printed pencil in (4.2.3)

Alternative : a proposed replacement pencil with an isolated ordinary quadratic singularity

Conserver le pinceau fourni dans (4.2.3). Le registre signale, pour n=2k+1 au moins 7, un plan singulier au lieu d’un point double ordinaire ; il distingue explicitement une réparation mathématique d’une lecture récupérée. La formule de remplacement détaillée n’est pas jointe et ne peut être inventée.

Statut : retained; proposed mathematical replacement unavailable. Source : \nolinkurl{conventions/GLOBAL_ISSUES_SGA.md:45}

\subsection*{ED-G48-inner-underlines — 7II:248 and other internal commas}

Lecture de la source conservée : source-layer continuous underlines through some inner commas

Alternative : stop and restart each underline around the comma

Conserver le soulignement fourni à chaque occurrence. Le registre précise que le dactylographe passait parfois sous les virgules et parfois non ; aucune substitution globale n’est justifiée sans le scan de chaque lieu.

Statut : needs the scan. Source : \nolinkurl{conventions/GLOBAL_ISSUES_SGA.md:48}

\subsection*{ED-G58-comma — 7II:332}

Lecture de la source conservée : H\_s,

Alternative : H\_\{s,\}

Conserver la virgule comme séparation des arguments, après H\_s. La position basse observée à faible confiance ne définit pas un nouvel indice. Le scan est nécessaire pour examiner la ligne de base physique.

Statut : needs the scan. Source : \nolinkurl{conventions/GLOBAL_ISSUES_SGA.md:58}

\subsection*{ED-G63-miracle — 7II:416}

Lecture de la source conservée : oh miracle (both attached layers)

Alternative : ô miracle (earlier held reading)

Conserver oh miracle, lecture commune des deux couches fournies. L’observation antérieure ô miracle n’avait pas été établie ; seule la consultation du scan permettrait de trancher indépendamment.

Statut : needs the scan. Source : \nolinkurl{conventions/GLOBAL_ISSUES_SGA.md:63}

\subsection*{ED-G64-quotient — 7II:423}

Lecture de la source conservée : the displayed fraction in (4.5.21)

Alternative : a slash presentation of the same quotient

Conserver la fraction de (4.5.21). Les deux couches donnent les mêmes numérateur et dénominateur ; la barre oblique proposée concerne uniquement la disposition graphique du quotient. Le scan est nécessaire pour établir cette disposition.

Statut : needs the scan. Source : \nolinkurl{conventions/GLOBAL_ISSUES_SGA.md:64}

\subsection*{ED-G64-subset-stroke — 7II:424}

Lecture de la source conservée : \textbackslash{}subseteq in the supplied source

Alternative : \textbackslash{}subset without the lower stroke

Conserver le signe subseteq de la couche fournie. Le trait inférieur peut relever d’une surfrappe ; la convention historique de subset peut aussi être non stricte. Le sens seul ne décide donc pas de la forme physique.

Statut : needs the scan. Source : \nolinkurl{conventions/GLOBAL_ISSUES_SGA.md:64}

\subsection*{ED-G64-Delta-Lambda — 7II:424–429}

Lecture de la source conservée : \textbackslash{}Delta\_B throughout the defining discussion; \textbackslash{}Lambda\_B in printed (5.0.12)

Alternative : change all seven \textbackslash{}Delta\_B to \textbackslash{}Lambda\_B, or retain current \textbackslash{}Delta\_B at (5.0.12)

Conserver les glyphes de la couche diplomatique et la correction locale Delta de la couche courante. Delta\_B est défini en (4.5.28) et la formule (5.0.12) renvoie à la même série caractéristique. Le remplacement global des sept occurrences demanderait leurs scans et n’est pas effectué.

Statut : needs the scan. Source : \nolinkurl{conventions/GLOBAL_ISSUES_SGA.md:64}

\subsection*{ED-G67-Dwork-psi — 7II:436–438}

Lecture de la source conservée : \textbackslash{}psi\_q in the attached layers

Alternative : \textbackslash{}Psi\_q suggested by the typewriter capital-form glyph

Conserver psi pour l’opérateur de Dwork déjà défini dans l’exposé. La hauteur de capitale du glyphe dactylographié ne suffit pas à établir un second opérateur ni une distinction de casse. La convention 78 exclut les simples variantes de forme du texte diplomatique.

Statut : needs the scan. Source : \nolinkurl{conventions/GLOBAL_ISSUES_SGA.md:67}

\subsection*{ED-G73-statement-underlines — 7II:403,405 and similar statements}

Lecture de la source conservée : source underlines sometimes include mathematical symbols and final punctuation

Alternative : split each underline around symbols and punctuation according to the print

Conserver le balisage fourni, sauf les réparations précises documentées séparément. Le registre renvoie ces soulignements à la vérification visuelle au cas par cas de la convention 48 ; aucune règle globale ne permet de retrouver le geste du dactylographe.

Statut : needs the scan. Source : \nolinkurl{conventions/GLOBAL_ISSUES_SGA.md:73}
